\section{Background and related work}

\paragraph{Solar resource physics.} The astronomical geometry of insolation ---
declination, hour angle, and the daily integral of extraterrestrial radiation
on a horizontal surface --- is textbook material \cite{duffie2013solar}. The
partition of incoming sunlight into reflection, atmospheric absorption, and
surface absorption is quantified by global energy-budget studies
\cite{trenberth2009earth}, and the accepted total solar irradiance is
$1361~\mathrm{W/m^2}$ \cite{kopp2011total}. We reuse these closed forms as
computable ground truth rather than fitting data.

\paragraph{Space-based solar power.} Collecting sunlight in orbit and beaming
it down was proposed by Glaser \cite{glaser1968power} and has been revisited in
modern phased-array concepts \cite{mankins2014sps}. It is the practical cousin
of Dyson's stellar-scale collector \cite{dyson1960search}: rather than
enclosing the star, place a collector where the light is undiminished and
uninterrupted. We treat it quantitatively as one option among several, not as
science fiction.

\paragraph{Sharing land and light with ecosystems.} Human appropriation of net
primary production already claims a large share of the biosphere's energy
throughput \cite{haberl2007quantifying}. Agrivoltaics shows that photovoltaics
and biological light demand can co-exist on the same land when light is shared
rather than monopolised \cite{barron2019agrivoltaics}. Our ecosystem-aware
optimisation abstracts that idea to a schedulable light budget and asks for the
allocation that maximises harvest while keeping each ecosystem whole.

\paragraph{Our contribution.} We unify the resource ceiling, land feasibility,
and the orbital alternative in one reproducible budget, and add a novel
optimisation layer: because an ecosystem's light requirement is a daily
integral, not an instantaneous one, the harvest schedule has temporal slack we
can exploit. We characterise the optimum in closed form and measure the energy
it unlocks.
